% #############################################################################
% This is Chapter 2
% !TEX root = ../main.tex
% #############################################################################
% Change the Name of the Chapter i the following line
\fancychapter{Background}
% The following line allows to ref this chapter
\label{chap:back}

DEPT combines three main ingredients: a bit-level representation of network traffic, a
heterogeneous detection ensemble, and a learned purifier. This chapter presents the
concepts each of them builds upon. Instead of a broad survey of machine learning for
security, the focus is kept on what the system relies on, so that the design decisions of
\Cref{chap:solution} follow from the approaches, techniques, and principles introduced here.

% #############################################################################
\section{Machine Learning for NIDS}\label{sec:bg-ml-nids}
A \ac{NIDS} inspects network traffic in order to separate benign activity from attacks.
Classical systems match traffic against hand-written signatures of known threats, which
makes them precise on catalogued attacks but unable to detect anything absent from the rule
set, including minor variants of existing malware. \ac{ML} reformulates detection as a
data-driven problem: instead of enumerating signatures, a model learns the statistical
structure that separates benign from malicious traffic, and can therefore generalize to
variants and to previously unseen attacks~\cite{boutaba2018comprehensive,wang2021machine}.
This shift supports a large body of modern detection
research~\cite{xu2025deep,anis2025survey}, and it defines the setting in which DEPT
operates.

\subsection{Learning Paradigms}\label{subsec:bg-paradigms}
Learning-based intrusion detection is usually organized by the kind of supervision a model
requires.

\emph{Supervised learning} trains on labelled data, where each example is marked benign or
malicious, and learns a mapping from input features to labels (\Cref{fig:bg-paradigms}).
Random forests and support vector machines follow this paradigm, as do deep neural networks
when trained as classifiers~\cite{boutaba2018comprehensive,sarker2021deep}; a neural
network is an architecture rather than a paradigm, however, and can equally be trained
without labels, as the autoencoders discussed below are. Every component of DEPT, the
two base detectors, the meta-learner that fuses them, and the purifier, is trained under
supervision. \emph{Unsupervised learning} uses no attack labels. It models
the normal behaviour of a network and flags departures from that baseline as anomalies,
which makes it attractive for catching zero-day attacks that a supervised model never saw
during training. A canonical example in this domain is Kitsune, an ensemble of
autoencoders that scores each packet by how poorly it reconstructs under a model of benign
traffic~\cite{mirsky2018kitsune}; it is also one of the independent detectors DEPT is
evaluated in front of in \Cref{chap:evaluation}. \emph{Semi-supervised learning} sits between the two,
exploiting a small labelled set together with abundant unlabelled data, and
\emph{reinforcement learning} replaces the fixed dataset with an agent that learns by
interacting with an environment. DEPT uses neither semi-supervised nor reinforcement
learning, although the latter is the natural framework for the continuous-adaptation extension discussed in
\Cref{chap:conclusion}.
\emph{Self-supervised learning} is a variant of unsupervised learning, deriving its
training signal from the data itself rather than from annotations: TANTGen~\cite{pissarra2026tantgen}, one of the key adversaries
DEPT aims to be protected against, is pre-trained to reconstruct masked packet sizes and inter-arrival
times on unlabelled traffic, and only its surrogate detector head is trained under
supervision~\cite{pissarra2026tantgen}.

\begin{figure}[htbp]
    \centering
    \subfigure[Supervised learning: labelled flows train a model that maps a flow to a label.]{\includesvg[width=0.46\textwidth]{./Images/supervised_learning}}
    \hfill
    \subfigure[Unsupervised learning: with no labels, the model learns the structure of the traffic and flags what departs from it.]{\includesvg[width=0.46\textwidth]{./Images/unsupervised_learning}}
    \caption[The two supervision regimes most relevant to DEPT.]{The two supervision regimes most relevant to DEPT, on network flows. All of DEPT's components are supervised.}
    \label{fig:bg-paradigms}
\end{figure}

\subsection{Feature Engineering and Representation}\label{subsec:bg-features}
The performance of a learned detector depends directly on the representation it consumes.
Feature engineering --- in this case, the process of turning raw packets into numerical vectors --- places an
upper bound on what any model can learn~\cite{salim2026feature,ning2024study}. Before
discussing representations, one term should be fixed: throughout this document, a
\emph{flow} is the sequence of packets that share the same five-tuple (source and
destination \ac{IP} addresses, source and destination ports, and transport protocol), ordered by
arrival time.

Two axes organize the design space of traffic representations. The first is
\emph{granularity}. \emph{Static}, packet-level features are read from individual headers
(addresses, ports, flags); they are cheap enough to extract at line rate, but they carry no
temporal context. \emph{Dynamic}, flow-level features aggregate a sequence of packets into
statistics such as mean packet size, inter-arrival times, or flow duration. These capture
behaviour, but at the cost of latency, since the system must wait for the flow to
progress~\cite{apruzzese2023sok}. The second axis concerns the space in which a
perturbation is expressed: the \emph{traffic-space} of actual packets on the wire versus
the \emph{feature-space} of derived vectors. This distinction becomes central once
adversaries are considered (\Cref{subsec:bg-realizability}). DEPT works at both
granularities, and deliberately keeps its features close to the traffic-space so that its
outputs can be turned back into real packets.

Rather than hand-designing features, DEPT represents traffic with \textbf{nPrint}, a
bit-level encoding of packet headers~\cite{holland2021nprint}. Hand-crafted feature sets
usually assume in advance which fields matter, and they are tied to particular protocols: a
tool that computes flow statistics for \ac{TCP} says nothing useful about an \ac{ICMP} flood. nPrint
instead maps each bit of a packet's headers to one feature in a fixed-width vector, using a
tri-state encoding: a value of $1$ or $0$ records a header bit that is set or unset, and
$-1$ marks a field that belongs to a protocol layer absent from that packet. The
representation is uniform across protocols, requires no manual feature design, and can be
consumed directly by models that learn their own features. Bit-level encodings of this kind
have proven effective for traffic analysis in settings where hand-crafted features
generalize poorly~\cite{holland2021nprint,wang2017end,lotfollahi2020deep}. \Cref{chap:solution}
details the exact encoding DEPT derives from it.

A bit-level header vector does, however, inevitably include fields that identify the
particular hosts and ports of a capture. If these are left intact, a model can reach high
benchmark accuracy simply by memorizing which hosts appear in the attack traces, instead
of learning the attack behaviour itself; this shortcut inflates test scores but does not
survive deployment on a different network~\cite{arp2022dos,apruzzese2023sok}, which is
why DEPT masks those fields (\Cref{chap:solution}).

\subsection{Model Architectures Relevant to DEPT}\label{subsec:bg-architectures}
DEPT combines three architecture families, chosen for complementary properties rather than
for raw accuracy alone.

\textbf{Random forests.}
An \ac{RF} is an ensemble of decision trees trained by bagging (\Cref{subsec:bg-ensemble-types}):
each tree recursively partitions the feature space with axis-aligned threshold tests, and
the forest averages their votes~\cite{breiman2001random,breiman1996bagging}. Forests are a
strong, low-cost fit for tabular and statistical features and need no feature
scaling~\cite{lucas2023comprehensive}, which is why DEPT applies one to flow-level
statistics. Their decisive property for this work, however, is
structural: a forest's decision function is constant over axis-aligned regions and
therefore \emph{non-differentiable}. It exposes no informative gradient that an attacker
can follow, a property whose relevance for adversarial robustness is developed in
\Cref{subsec:bg-gradient-masking}.

\textbf{One-dimensional and temporal convolutional networks.}
A \ac{CNN} slides learnable kernels over its input to detect local patterns; in one
dimension it processes sequences, extracting recurring motifs from a packet stream such as
header patterns or sizing signatures~\cite{wang2017end,lotfollahi2020deep}
(\Cref{fig:bg-1dcnn}). A \ac{TCN} specializes this design for long sequences through two
devices~\cite{bai2018empirical}. \emph{Dilated} convolutions insert a growing stride
between kernel taps, so the receptive field expands with depth without a corresponding
blow-up in parameters~\cite{oord2016wavenet}. \emph{Causal} convolutions pad the input so
that the output at position $t$ depends only on positions $\leq t$, never on future
packets. Causality is what makes the model valid for streaming inference: it can emit a
verdict after each packet arrives.

\begin{figure}[htbp]
    \centering
    \includesvg[width=0.92\textwidth]{./Images/one_d_cnn}
    \caption[The one-dimensional convolutional pipeline.]{A one-dimensional convolutional pipeline over a packet sequence: kernels slide along the sequence and produce per-position feature maps, which further layers, pooling, and a classification head turn into a verdict. DEPT's temporal convolutional network specializes this design with dilated, causal convolutions.}
    \label{fig:bg-1dcnn}
\end{figure}

\textbf{Transformers.}
The self-attention mechanism lets every position in a sequence attend to every other
position, weighting them by a learned measure of
relevance~\cite{vaswani2017attention}. Each input position is projected into a query, a
key, and a value vector; collecting these as the rows of matrices $Q$, $K$, and $V$,
attention computes
\begin{equation}
    \mathrm{Attention}(Q, K, V) \;=\; \mathrm{softmax}\!\left(\frac{QK^{\top}}{\sqrt{d_k}}\right)V,
    \label{eq:bg-attention}
\end{equation}
where $d_k$ is the key dimension, so the output at each position is a weighted combination
of the values of the positions it attends to. This captures long-range dependencies that a
fixed convolution window cannot. Two arrangements are relevant here. \emph{Encoder-only}
models attend bidirectionally over the entire sequence, which suits classification and
representation tasks, including learned traffic
analysis~\cite{lin2022bert,ghadermazi2025gtae}. \emph{Decoder-only} models apply a causal
mask so that each position attends only to the past (\Cref{fig:bg-transformer}); this is
the arrangement behind autoregressive generation, but the mask is needed by any task whose
output at a position may depend only on that position's past, generative or not.
Causality also enables \emph{key--value caching}: the keys and values of past positions in
\Cref{eq:bg-attention} do not change as new tokens arrive, so they can be stored and reused
instead of reprocessing the whole prefix on every packet.

\begin{figure}[htbp]
    \centering
    \includesvg[width=0.42\textwidth]{./Images/decoder_only_transformer}
    \caption[The decoder-only Transformer block.]{A decoder-only Transformer block over a stream of packet tokens: masked self-attention, then a feed-forward sublayer, with residual connections and layer normalization, stacked $N$ times. The causal mask restricts each position to its past, which is what makes the design valid for streaming inference. Both of DEPT's Transformer components use this block.}
    \label{fig:bg-transformer}
\end{figure}

% #############################################################################
\section{Adversarial Machine Learning}\label{sec:bg-aml}
The same flexibility that lets a learned model generalize also exposes it to manipulation.
An \emph{adversarial example} is an input perturbed by a small, deliberate amount so that
the model misclassifies it, often while remaining nearly indistinguishable from a clean
input~\cite{szegedy2014intriguing,goodfellow2014explaining}. For a \ac{NIDS}, adversarial
examples are the mechanism of \emph{evasion}: an attacker reshapes malicious traffic just
enough for it to be labelled benign~\cite{he2023adversarial,ennaji2025adversarial}.
\ac{AML} studies both how such manipulations are constructed and how models can be hardened
against them.

\subsection{Threat Taxonomy}\label{subsec:bg-threat-taxonomy}
Attacks are most commonly classified along two axes: what the adversary knows, and what the adversary
wants.

Along the knowledge axis, a \emph{white-box} attacker has full access to the model,
including its architecture and weights, and can compute gradients of the loss with respect
to the input to craft precise perturbations. A \emph{gray-box} attacker knows only part of
the pipeline, such as the feature set. A \emph{black-box} attacker can only query the
system and observe its outputs. The black-box setting is the realistic one for a deployed
\ac{NIDS}, which an attacker probes as an oracle, and such attacks typically rely on
\emph{transferability}: an adversarial example crafted against a substitute model often
fools the target as well~\cite{papernot2017practical}. Along the goal axis, \emph{poisoning}
corrupts the model during training, whereas \emph{evasion} manipulates inputs at inference
time against an already-trained model. DEPT is designed for the black-box evasion threat,
which is the case that dominates real deployments (\Cref{chap:solution}).

Network intrusion detection differs from the image classification setting, where
adversarial examples were first studied, by a constraint with no equivalent in the pixel
domain: a perturbed feature vector must correspond to a packet or flow that is actually
valid on the wire. Recovering a realizable packet from a desired feature vector is the
\textbf{\ac{IFM}} problem, and it is frequently constrained or computationally
intractable~\cite{pierazzi2020intriguing,hashemi2019towards,sun2023gpmt};
\Cref{subsec:bg-realizability} returns to it. In particular, the gradient-based methods
that dominate the image setting, such as the \ac{FGSM} and
\ac{PGD}~\cite{goodfellow2014explaining,madry2018towards}, produce feature perturbations
that need not correspond to any sendable packet, which limits their direct use in this
domain.

\subsection{Defense Paradigms}\label{subsec:bg-defense-paradigms}
Defenses against evasion fall into a few broad families, with different strengths for
traffic data. \emph{Adversarial training} augments the training set with adversarial
examples so that the model learns to classify them correctly~\cite{madry2018towards}. It is
effective against the perturbations it is trained on, but it is costly, it can reduce
accuracy on clean traffic, and in networking it is only meaningful if the adversarial
examples are themselves realizable rather than impossible feature vectors. DEPT's
detectors are likewise trained on realizable adversarial traffic, although with a
different goal: rather than learning to keep classifying a perturbed attack as an attack,
they learn to recognize that a flow has been manipulated at all (\Cref{sec:rel-training}).

\emph{Input transformation}, or \emph{purification}, is the family DEPT builds its purifier
on. Instead of hardening the classifier, it transforms each input before classification in
order to remove the perturbation, projecting the sample back onto the manifold of normal
traffic, as \Cref{fig:bg-purification} illustrates. The idea originates with Defense-GAN,
which projects each input onto the manifold captured by a generative adversarial
network~\cite{samangouei2018defense}, and it has since moved to diffusion
models~\cite{li2025adbm} and to the network domain~\cite{alhussien2024advpurrec,merzouk2024diffusion}.
DEPT's purifier pursues the same
goal without being a generative model itself: it is a causal Transformer denoiser,
trained on matched adversarial-to-clean pairs rather than on a model of the benign
distribution it could sample from (\Cref{chap:solution}).

\begin{figure}[htbp]
    \centering
    \includesvg[width=0.80\textwidth]{./Images/purification_manifold}
    \caption[Purification as projection onto the benign manifold.]{Purification as projection onto the benign manifold. A perturbation $\delta$ moves a clean flow $x$ off the manifold, and the detector misclassifies the result $x'$. The purifier $g$ projects $x'$ back, so the detector receives an input close to $x$. The idea originates with Defense-GAN~\cite{samangouei2018defense}.}
    \label{fig:bg-purification}
\end{figure}

DEPT adopts neither of the two remaining families. \emph{Defensive distillation} smooths
the decision surface so that useful gradients become harder to
find~\cite{papernot2016distillation}. \emph{Certified robustness} guarantees that no perturbation within a
bounded region changes a prediction, through randomized
smoothing~\cite{cohen2019certified,yan2025certta} or Lipschitz and edit-distance
bounds~\cite{abad2025lipslev}; such guarantees are hard to scale to the high-dimensional,
discrete structure of traffic, and they secure a formal radius rather than the
traffic-space validity that is central to this work.

\subsection{Realizability and Protocol Constraints}\label{subsec:bg-realizability}
The gap between a \emph{valid feature vector} and a \emph{valid packet} is a central
difficulty of network adversarial machine learning. A perturbation can move a feature
vector anywhere in $\mathbb{R}^{n}$, but a packet must satisfy hard structural constraints:
checksums must match the bytes they cover, header-length fields must be equal to the actual
lengths, control flags must be mutually consistent, and sequence numbers must be ordered. A
perturbation that ignores these constraints is \emph{unrealizable}: it describes no packet
that can be sent, so evaluating a detector against it overestimates the real
threat~\cite{pierazzi2020intriguing,hashemi2019towards}. Recent attack research enforces
realizability explicitly, whether by using an SMT solver to find packets that satisfy
the protocol constraints, as in PANTS~\cite{jin2025robustifying}, or by generating
traffic that is protocol-compliant by design, as in GPMT~\cite{sun2023gpmt} and
TANTGen~\cite{pissarra2026tantgen}.

The constraint applies to both sides, which is what makes it a design concern for DEPT
rather than a property of attacks alone. A defense that \emph{edits} traffic, as a purifier
does, must itself re-emit valid packets; otherwise its output is unusable, or is trivially
detected as malformed. The step from feature-space back to traffic-space, which we call
\emph{realization}, is therefore treated as a first-class part of the design. The nPrint
representation helps here: because it is a bit-level image of the packet headers, a
purified nPrint vector can be decoded back into a concrete packet and checked for
on-the-wire validity, which DEPT does with packet-level tools (\Cref{chap:solution}).

% #############################################################################
\section{Ensemble Learning for Robustness}\label{sec:bg-ensembles}
An \emph{ensemble} combines several base models into a single predictor, on the premise
that a collective decision is more accurate and more robust than that of any single
learner~\cite{lucas2023comprehensive,bibers2025ensemble}. Beyond accuracy, ensembles are of
particular interest for security, because diversity among the base models can be turned
directly against adversarial transferability.

\subsection{Bagging, Boosting, Stacking}\label{subsec:bg-ensemble-types}
Three combination strategies dominate (\Cref{fig:bg-ensembles}). \emph{Bagging} trains
homogeneous models in parallel on bootstrap resamples of the data and averages their
predictions, which reduces variance; the random forest is bagging applied to decision
trees~\cite{breiman1996bagging,breiman2001random}. \emph{Boosting} trains models
sequentially, each one focusing on the examples its predecessors got wrong, which reduces
bias~\cite{freund1997decision}. \emph{Stacking} trains \emph{heterogeneous} base learners
and then trains a separate \emph{meta-learner} to combine their
outputs~\cite{wolpert1992stacked}. Stacking is most effective when the base
learners are diverse and make different kinds of errors, so that their combination covers
more of the input space than any of them alone.

\begin{figure}[htbp]
    \centering
    \subfigure[Bagging: parallel homogeneous models on resampled data.]{\includesvg[width=0.31\textwidth]{./Images/bagging}}
    \hfill
    \subfigure[Boosting: sequential models, each correcting the last.]{\includesvg[width=0.31\textwidth]{./Images/boosting}}
    \hfill
    \subfigure[Stacking: heterogeneous models fused by a meta-learner.]{\includesvg[width=0.31\textwidth]{./Images/stacking}}
    \caption[The three classical ensemble strategies.]{The three classical ensemble strategies. DEPT's detector is a stacking ensemble, chosen because its base learners are architecturally heterogeneous and emit predictions at different rates.}
    \label{fig:bg-ensembles}
\end{figure}

Stacking is also the natural fit when the base learners are heterogeneous not only in
architecture but in operation, for example when one emits a verdict on every packet and
another only at sparse points of a flow. A fixed combination rule, such as averaging or
majority voting, cannot express when each learner should be trusted. A learned
meta-learner can: by assigning input and position-dependent weights to each base output,
it acts as a dynamic gating mechanism rather than a static
average~\cite{yang2025novel,gungor2022stewart}.

\subsection{Gradient Masking and Transferability}\label{subsec:bg-gradient-masking}
DEPT is itself a learned detector, so a sophisticated attacker may also try to evade DEPT
directly. Such an attacker typically knows neither which models DEPT uses nor their
parameters, the black-box setting of \Cref{subsec:bg-threat-taxonomy}.
Transferability is what makes black-box evasion practical: an attacker who cannot see the
target model builds a differentiable substitute, crafts a perturbation against it with a
gradient method such as \ac{PGD} or \ac{FGSM}, and relies on that perturbation carrying
over to the target~\cite{papernot2017practical}. Heterogeneous stacking is a way of
increasing DEPT's own robustness against this mechanism. A differentiable member, such as the \ac{TCN}, exposes a gradient that an
attacker can follow. An \ac{RF} does not: its decision function is constant over
axis-aligned regions, so there is no gradient direction to exploit. When a differentiable
and a non-differentiable model are combined, a perturbation optimized against the
differentiable member does not necessarily move the forest's discrete verdict, so it
transfers poorly across the ensemble and raises the attacker's
cost~\cite{yang2021trs,gungor2022stewart,biggio2011bagging,chen2022collective}.

This argument must be stated carefully, because the absence of a usable gradient is a form
of \emph{gradient masking}, and gradient masking is not robustness in itself. A
determined white-box adversary can estimate gradients numerically or resort to
gradient-free search, and masked gradients can produce a false sense of
security when mistaken for genuine robustness~\cite{athalye2018obfuscated}. The argument
holds for DEPT precisely because the attackers it is designed against have no white-box
access to DEPT as a whole: the value of a
non-differentiable member is not an absolute barrier, but a measurable reduction in the
transferability from whatever substitute the attacker is able to build without direct
gradient access. Robustness is therefore treated as a property to be evaluated empirically,
not as a guarantee to be assumed, which is what \Cref{chap:evaluation} does.

% #############################################################################
\section{Summary}\label{sec:bg-summary}
The concepts in this chapter compose into the three pillars of DEPT.
\emph{Representation}: a bit-level nPrint encoding stripped of identity fields, paired with
prefix statistics measured on the same packets, gives heterogeneous learners complementary
behaviour-focused and realizable views of traffic (\Cref{subsec:bg-features}). \emph{Purification}: an input-transformation defense, built as a
causal Transformer, projects flagged traffic back onto the benign manifold while respecting
the protocol constraints that keep its output realizable
(\Cref{subsec:bg-defense-paradigms,subsec:bg-realizability}). \emph{Ensembling}: heterogeneous stacking, fused by
an attention meta-learner, combines a differentiable network with a non-differentiable
forest to resist transferable gradient attacks under a black-box threat model
(\Cref{sec:bg-ensembles}). 
